%!TEX root = ../graduate-work.tex
\section{Applying Parallelism to Neural Network Verification}
\label{sec:parallel-verifiers}

This chapter justifies why Tensor Parallelism and FSDP are suited to bound
propagation, analyses their accuracy--memory trade-offs, and discusses
integration with complete verification.

\subsection{Why Tensor Parallelism Is Applicable to Verification}

Bound-propagation algorithms (CROWN, $\alpha$-CROWN) reduce bound computation
to a chain of matrix multiplications executed as a backward pass over the
network's computational graph. For a linear layer with weight matrix $W$, the
CROWN backward performs $A \leftarrow A \cdot W$, where $A$ is the matrix of
linear-relaxation coefficients. This operation is identical to the backward
pass during training, so TP transfers almost without modification.

For fully connected networks, weight matrices dominate GPU memory; TP shards
exactly these, achieving close-to-linear reduction. The communication pattern
also matches: one \texttt{AllReduce} per Column--Row pair, with all intermediate
computations remaining local. Crucially, TP does not alter the topology of the
computational graph: PP, by contrast, partitions the model into stages, and
CROWN-backward then cannot traverse all layers in a single pass---inadmissible
for verification. \texttt{auto\_LiRPA} builds the bound-propagation graph on
top of the original network, and a single connected graph is a necessary
condition.

\subsection{Impact of Tensor Parallelism on Bound Tightness}
\label{sec:tp-precision-loss}

The key drawback of TP is loss of bound tightness, which grows with the number
of sharded zones. This loss is not an implementation artifact: it is
fundamentally caused by the interaction of bound propagation with partial weight
matrices.

\subsubsection{Backward Pass in CROWN}

CROWN computes linear bounds on the network output as a function of the input.
For a network with $d$ layers and weight matrices $W_1, \ldots, W_d$, the
backward pass constructs the $A$-matrix---the product of linearised weights:
$$A = \Lambda_d W_d \cdot \Lambda_{d-1} W_{d-1} \cdots \Lambda_1 W_1,$$
where $\Lambda_i = \mathrm{diag}(\lambda_i)$ is the diagonal matrix of
ReLU linear-relaxation slopes (values in $[0, 1]$, determined from intermediate
neuron bounds).

The resulting bounds take the form:
$$\ell \leq A \cdot \mathbf{x} + \mathbf{b} \leq u,$$
where $\mathbf{x}$ is the input vector within an $\varepsilon$-ball.
The key property: $A$ is a single linear function of the \textit{original}
input, preserving dependencies between variables across all layers.

\subsubsection{The Challenge of CROWN under TP}

Under TP, each Column--Row pair forms a \textit{sharded zone}---a computational
subgraph in which activations have dimension $h/P$ per GPU instead of~$h$.

Let layer $W_i \in \mathbb{R}^{h \times h}$ be column-sharded: GPU~$r$ stores
$W_i^{(r)} \in \mathbb{R}^{(h/P) \times h}$. During CROWN backward, the
$A$-matrix from the preceding layer has full dimension~$h$:
$$A \leftarrow A \cdot \Lambda_i W_i.$$
Each GPU needs the \textit{full} matrix $W_i$ for this multiplication, yet
each holds only $W_i^{(r)}$---$h/P$ rows---making a correct multiplication
without communication impossible.

For a Column--Row pair with no intermediate layers, the problem is resolved by
a simple \texttt{AllReduce}, which recovers the full result. However, when ReLU
activations lie between the Column and Row layers (as in typical networks with
more than two layers), the situation becomes more complex: to select
relaxation slopes $\lambda_i$, CROWN requires \textit{intermediate bounds}
of neurons in the sharded zone, and computing them exactly would require a
separate backward pass with communication at each step---negating the memory
savings and rendering TP impractical for deep networks. The only practical
alternative is to use IBP for intermediate bounds within zones, since IBP
operates layer-wise with local weights and requires no $A$-matrices.

\subsubsection{Why IBP Is Substantially Less Tight}

The difference between IBP and CROWN lies in how each method handles
dependencies between variables.

\textbf{CROWN} maintains a global linear dependency
$y = A \mathbf{x} + \mathbf{b}$, where $A$ encodes the influence of each input
component through \textit{all} intermediate layers. Bounds are computed once:
$$lb = \sum_j \min(A_j x_j^L, A_j x_j^U) + b, \quad
  ub = \sum_j \max(A_j x_j^L, A_j x_j^U) + b.$$

\textbf{IBP} computes bounds layer by layer. For the $i$-th layer with weights
$W_i$ and input interval $[\ell_{i-1}, u_{i-1}]$:
\begin{align*}
\ell_i &= W_i^+ \ell_{i-1} + W_i^- u_{i-1} + b_i, \\
u_i &= W_i^+ u_{i-1} + W_i^- \ell_{i-1} + b_i,
\end{align*}
where $W_i^+ = \max(W_i, 0)$, $W_i^- = \min(W_i, 0)$.

The weakness of IBP is componentwise interval treatment: variable dependencies
are not tracked. This gives rise to the \textit{wrapping effect}~\cite{moore2009interval}---an
artifact of interval arithmetic in which bound widths can grow exponentially
with network depth.

\textbf{Example.} Consider a two-dimensional input
$\mathbf{x} = (x_1, x_2)^T$ with $x_1, x_2 \in [-1, 1]$ and
$W = \begin{pmatrix} 1 & 1 \\ 1 & -1 \end{pmatrix}$.

CROWN: $y = W\mathbf{x}$, so $y_1 = x_1 + x_2 \in [-2, 2]$,
$y_2 = x_1 - x_2 \in [-2, 2]$. Applying $W$ again:
$z_1 = y_1 + y_2 = 2x_1 \in [-2, 2]$.

IBP: the first layer gives $y_1 \in [-2, 2]$, $y_2 \in [-2, 2]$---matching
CROWN. But when applying $W$ again, IBP treats $y_1$ and $y_2$ as
\textit{independent}: $z_1 = y_1 + y_2 \in [-4, 4]$---twice as wide as the
true bound. CROWN, by tracking the dependencies $y_1 = x_1 + x_2$,
$y_2 = x_1 - x_2$, yields the exact result.

\subsubsection{Accumulation of Error with Depth}

For a random weight matrix with entries of order $O(1/\sqrt{h})$ (typical
initialisation), the IBP interval width after $k$ layers grows as $O(\rho^k)$,
where $\rho$ is the spectral radius of the generalised absolute-value matrix $|W|$.

The experimental data (Table~\ref{tab:tp-correctness}) confirm this dependence.
With one Column--Row pair (model 256$\times$2), the sharded zone contains no
intermediate ReLU; IBP fallback is not used, and the deviation is
$\sim\!10^{-8}$ (float32 machine precision). With two pairs (256$\times$4),
each zone contains an intermediate ReLU, IBP is used, and the deviation is
$\sim\!2.5$---an increase of eight orders of magnitude. With three pairs
(256$\times$6), three zones each containing a ReLU yield a deviation of
$\sim\!750$.

The root cause is a \textit{method substitution}, not a distributed-computation
error: IBP must be used instead of CROWN inside sharded zones, and with each
additional zone the error compounds over depth.

TP-based verification is effective for shallow models (1--2 sharded zones).
For deep models ($\geq\!4$ zones), the loss of bound tightness renders
bounds practically useless---motivating the FSDP approach
(Section~\ref{sec:fsdp-memory-theory}), which has no such limitation.

\subsection{Comparison of TP and FSDP for Verification}

Pipeline Parallelism is inapplicable to verification: CROWN performs a backward
pass over \textit{all} layers, and with PP, where each GPU stores only its
stage, a complete backward pass would require shipping weights between GPUs---
effectively turning PP into a less efficient variant of TP with point-to-point
communication.

FSDP shards weights across GPUs, assembling them via \texttt{AllGather}
immediately before each layer. In the verification context this is a fundamental
advantage over TP: each layer receives the full weight matrix, so computations
are mathematically \textit{equivalent} to the single-GPU case without any loss
of bound tightness. The cost is higher communication overhead:
\texttt{AllGather} before every layer instead of one \texttt{AllReduce} per
layer pair in TP.

TP shards both weights and $A$-matrices, achieving up to ${\sim}2\times$ peak
memory reduction but with degrading bound tightness as the number of sharded
zones grows. FSDP shards \textit{only weights}, achieving exactly 50\% baseline
memory savings (factor $1/P$) and 17--39\% peak memory savings, with bounds
bitwise identical to the single-GPU baseline---critical for complete verification
($\beta$-CROWN and Branch-and-Bound).

\subsection{TP Implementation for Bound Propagation in \texttt{auto\_LiRPA}}

Integration of TP into \texttt{auto\_LiRPA}~\cite{AutoLiRPARepo} required four
components.

Operators \texttt{BoundLinearTP\_Col} and \texttt{BoundLinearTP\_Row} inherit
from \texttt{BoundLinear}, store sharded weight matrices, and add
\texttt{AllReduce} to \texttt{bound\_backward}, ensuring correct bound
propagation through distributed layers.

Each operator is registered as a custom ONNX operation via
\texttt{torch.autograd.Function} and the \texttt{symbolic} method, allowing
the PyTorch JIT tracer to construct the computational graph correctly.

The function \texttt{tp\_shard\_bounded\_module} takes an arbitrary
\texttt{BoundedModule} (including one loaded from ONNX) and automatically
replaces alternating pairs of \texttt{BoundLinear} nodes with
\texttt{BoundLinearTP\_Col/Row}, sharding the weight matrices over the
available GPUs.

Between each Column--Row pair of nodes, a \textit{sharded zone} is formed---a
subgraph with reduced activation dimensionality. Nodes within zones (including
ReLU) are processed locally without GPU-to-GPU synchronisation; zone membership
of each node is determined by BFS labelling at module initialisation.

\subsection{FSDP Implementation for Bound Propagation in \texttt{auto\_LiRPA}}

The FSDP implementation requires no new operator types and does not modify the
computational graph---the total code amounts to approximately 100 lines, versus
${\sim}500$ for TP.

The function \texttt{fsdp\_shard\_bounded\_module} traverses the
\texttt{BoundedModule} graph, locates \texttt{BoundParams} nodes of fully
connected (\texttt{BoundLinear}) and convolutional (\texttt{BoundConv}) layers,
and replaces each matrix with a shard along dimension 0 (\texttt{dim=0}).
For \texttt{BoundLinear}, the matrix $W \in \mathbb{R}^{k \times n}$ is
replaced with $W^{(r)} \in \mathbb{R}^{(k/P) \times n}$; for
\texttt{BoundConv}, the tensor
$W \in \mathbb{R}^{C_\text{out} \times C_\text{in} \times k_H \times k_W}$
is replaced with
$W^{(r)} \in \mathbb{R}^{(C_\text{out}/P) \times C_\text{in} \times k_H \times k_W}$.
Sharding along the output dimension is correct since both layer types are
linear in the output and \texttt{AllGather} reconstructs the full matrix correctly.

When the \texttt{forward()} method of \texttt{BoundParams} detects the sharding
flag, it initiates \texttt{dist.all\_gather}, assembling the full matrix from
shards on all GPUs. The full matrix is returned as the forward-pass result and
used by all downstream operations without modification.

Hooks \texttt{fsdp\_gather\_node} and \texttt{fsdp\_free\_node} are integrated
into \texttt{backward\_general} (CROWN backward) and \texttt{IBP\_general}
(IBP forward): the full matrix is assembled immediately before the layer is
processed and freed immediately afterwards. At any given moment, at most one
full weight matrix is present in GPU memory (in addition to the shards).

\subsubsection{Theoretical Memory Analysis}
\label{sec:fsdp-memory-theory}

Let the model contain $d$ linear layers with weight matrices
$W_i \in \mathbb{R}^{h \times h}$, verification is performed by CROWN for a
single example, and $P$ is the number of GPUs.

In single-GPU mode all weights are stored in full: $M_W = d \cdot h^2$.
With FSDP, each GPU stores $1/P$ of each weight plus one full matrix for the
current layer:
$$M_W^{\text{FSDP}} = \frac{d \cdot h^2}{P} + h^2.$$

In CROWN backward, the $A$-matrix for intermediate layer $i$ has size
$O(h \times h)$. Under TP, $A$-matrices pass through sharded layers at reduced
dimension $h/P$; under FSDP they always have full dimension~$h$:
$$M_A^{\text{TP}} \propto \frac{h^2}{P}, \qquad
  M_A^{\text{FSDP}} \propto h^2.$$

Total peak memory for TP:
$$M_{\text{peak}}^{\text{TP}} \approx
    \frac{(d + C) \cdot h^2}{P},$$
yielding a linear reduction by a factor of $P$. For FSDP:
$$M_{\text{peak}}^{\text{FSDP}} \approx
    h^2 \!\left(\frac{d}{P} + 1 + C\right),$$
where $C$ is the number of simultaneously stored $A$-matrices.

FSDP saves memory only on weights, whereas TP reduces both terms. The ratio
of FSDP peak memory to single-GPU peak:
$$\frac{M_{\text{peak}}^{\text{FSDP}}}{M_{\text{peak}}^{\text{single}}}
    = \frac{d/P + 1 + C}{d + C}.$$
At $d = 4$, $C \approx 4$, $P = 2$ this gives
$\approx (2 + 1 + 4)/(4 + 4) = 7/8 = 0.875$---consistent with the
experimentally observed savings of ${\sim}34$\%.

Baseline memory (weight storage without $A$-matrices, measured before
\texttt{compute\_bounds}) is reduced exactly by a factor of $P$:
$$\frac{M_W^{\text{FSDP}}}{M_W^{\text{single}}} = \frac{1}{P},$$
since each GPU stores exactly $1/P$ rows of every weight tensor.
At $P = 2$ this gives $0.5$, matching the experimentally observed 50\%
reduction across all tested configurations.

\subsection{Challenges of Integrating TP with Complete Verification}

Integrating TP with BaB is significantly more complex than FSDP. Branch-and-Bound
selects a neuron to fix at each step; under TP, neurons in sharded layers are
distributed across GPUs, and even selecting the neuron requires collecting
statistics from all devices. Beyond this, $\beta$-CROWN encodes branching
constraints via dual variables $\beta$ that are defined locally for neurons in
sharded zones but affect global bounds through \texttt{AllReduce}---a non-trivial
synchronisation problem. Finally, BaB generates millions of sub-problems, and
each sub-problem in a batch must use the same TP configuration, complicating
load balancing.

FSDP is free of these difficulties: each rank receives full weight matrices via
\texttt{AllGather} and performs branching logic identical to the single-GPU case.
FSDP integration with $\beta$-CROWN + BaB is implemented and experimentally
confirmed (see Section~\ref{sec:exp6}): under an identical workload, both modes---
single-GPU and FSDP$=$2---produce bitwise-identical results, with \texttt{AllGather}
overhead not exceeding 5~MB.
